\documentclass[../../../include/open-logic-section]{subfiles}

\begin{document}

\olfileid{sth}{cardinals}{classing}
\olsection{المتناهي، و\usetoken{S}{enumerable}، و\usetoken{S}{nonenumerable}}

بعد أن تعرفنا إلى الكاردينالات، يجدر بنا أن نقضي بعض الوقت في الحديث عن
أنواعها المختلفة؛ وبالتحديد الكاردينالات المتناهية، و!!{enumerable}،
و!!{nonenumerable}.

تستلزم نتيجتانا الأوليان أن الكاردينالات المتناهية هي بالضبط الأعداد
الترتيبية المتناهية، التي عرّفناها من قبل بأنها \emph{أعدادنا الطبيعية}
في \olref[z][infinity-again]{defnomega}:

\begin{prop}\ollabel{finitecardisoequal}
لتكن $n, m \in \omega$. عندئذ $n = m$ إذا وفقط إذا
$\cardeq{n}{m}$.
\end{prop}

\begin{proof}
الاتجاه \emph{من اليسار إلى اليمين} بديهي. ولإثبات الاتجاه
\emph{من اليمين إلى اليسار}، افترض أن $\cardeq{n}{m}$ مع أن
$n \neq m$. وبحسب قانون التثليث، إما أن $n \in m$ وإما أن
$m \in n$؛ ولنفترض، من دون إخلال بالعموم، أن $n \in m$. عندئذ
$n \subsetneq m$ وتوجد !!a{bijection} $f \colon m \to n$، فيكون
$m$ غير متناهٍ بحسب ديدكند، خلافاً لـ
\olref[z][infinity-again]{naturalnumbersarentinfinite}.
\end{proof}

\begin{cor}\ollabel{naturalsarecardinals}
إذا كان $n \in \omega$، فإن $n$ كاردينال.
\end{cor}

\begin{proof}
وهذا مباشر.
\end{proof}
\noindent
وينتج أيضاً أن عدة تصورات معقولة لمعنى وصف كاردينال بأنه «متناهٍ» أو
«غير متناهٍ» تتطابق:
\begin{thm}\ollabel{generalinfinitycharacter}لكل مجموعة $A$، الشروط الآتية متكافئة:
\begin{enumerate}
	\item\ollabel{card:notinomega} $\card{A} \notin \omega$، أي إن
	$\card{A}$ ليس عدداً طبيعياً؛
	\item\ollabel{card:omegaplus} $\omega \leq \card{A}$؛
	\item\ollabel{card:infinite} $A$ غير متناهية بحسب ديدكند.
\end{enumerate}
\end{thm}

\begin{proof}
ينتج ذلك من \olref[ord-arithmetic][using-addition]{ordinfinitycharacter}
و\olref[cardinals][cardsasords]{lem:CardinalsBehaveRight}
و\olref{naturalsarecardinals}.
\end{proof}

وهذا يجيز لنا \emph{تعريف} بعض المفاهيم التي استعملناها استعمالاً غير
رسمي نوعاً ما في \olref[sfr][][]{part}:

\begin{defn}\ollabel{defnfinite}
نقول إن $A$ \emph{متناهية} إذا وفقط إذا كان $\card{A}$ عدداً طبيعياً،
أي إذا كان $\card{A} \in \omega$. وإلا قلنا إن $A$
\emph{غير متناهية}.
\end{defn}
\noindent
لكن لاحظ أن هذا التعريف يُعرض على خلفية $\ZFC$. فقد احتجنا، في نهاية
الأمر، إلى الترتيب الجيد كي نضمن أن لكل مجموعة كاردينالية. بل إنه من دون
الترتيب الجيد قد توجد مجموعة ليست متناهية ولا غير متناهية بحسب ديدكند.
سنعود إلى هذا النوع من المسائل في \olref[choice][]{chap}. أما الآن
فنواصل الاعتماد على الترتيب الجيد.

لننتقل الآن من الكاردينالات المتناهية إلى غير المتناهية. وفيما يأتي
نقطتان أوليتان:

\begin{cor}\ollabel{omegaisacardinal}
$\omega$ هو أصغر كاردينال غير متناهٍ.
\end{cor}

\begin{proof}
$\omega$ كاردينال، لأنه غير متناهٍ بحسب ديدكند، ولو كان
$\cardeq{\omega}{n}$ لأي $n \in \omega$ لكان $n$ غير متناهٍ بحسب
ديدكند، خلافاً لـ
\olref[z][infinity-again]{naturalnumbersarentinfinite}. أما كون
$\omega$ أصغر كاردينال غير متناهٍ فينتج من التعريف.
\end{proof}

\begin{cor}
كل كاردينال غير متناهٍ عدد ترتيبي حدي.
\end{cor}

\begin{proof}
ليكن $\alpha$ عدداً ترتيبياً خلفياً غير متناهٍ، بحيث
$\alpha = \beta \ordplus 1$ لبعض $\beta$. بحسب
\olref{finitecardisoequal} يكون $\beta$ غير متناهٍ أيضاً، ومن ثم
$\cardeq{\beta}{\beta \ordplus 1}$ بحسب
\olref[ord-arithmetic][using-addition]{ordinfinitycharacter}. والآن
$\card{\beta} = \card{\beta\ordplus 1} = \card{\alpha}$ بحسب
\olref[cardinals][cardsasords]{lem:CardinalsBehaveRight}، ولذلك
$\alpha \neq \card{\alpha}$.
\end{proof}

سبق أن أشرنا في \olref[sfr][siz][enm-alt]{defn:enumerable} إلى أننا
نستطيع التمييز بين المجموعات غير المتناهية !!{enumerable}
و!!{nonenumerable}. ويقود ذلك التعريف بصورة طبيعية إلى ما يأتي:

\begin{prop}
$A$ هي !!{enumerable} إذا وفقط إذا كان $\card{A} \leq \omega$، و$A$
هي !!{nonenumerable} إذا وفقط إذا كان $\omega < \card{A}$.
\end{prop}

\begin{proof}
بحسب قانون التثليث، الادعاءان متكافئان، فيكفي أن نثبت أن $A$ هي
!!{enumerable} إذا وفقط إذا كان $\card{A} \leq \omega$. في الاتجاه
\emph{من اليمين إلى اليسار}: إذا كان $\card{A} \leq \omega$، فإن
$\cardle{A}{\omega}$ بحسب
\olref[cardinals][cardsasords]{lem:CardinalsBehaveRight}
و\olref{omegaisacardinal}. وفي الاتجاه \emph{من اليسار إلى اليمين}:
افترض أن $A$ هي !!{enumerable}؛ فعندئذ، بحسب
\olref[sfr][siz][enm-alt]{defn:enumerable}، توجد ثلاث حالات ممكنة:
\begin{enumerate}
	\item إذا كان $A = \emptyset$، فإن $\card{A} = 0 \in \omega$، بحسب
	\olref{naturalsarecardinals} و
	\olref[cardinals][cardsasords]{lem:CardinalsBehaveRight}.
	\item إذا كان $\cardeq{n}{A}$، فإن $\card{A} = n \in \omega$، بحسب
	\olref{naturalsarecardinals} و
	\olref[cardinals][cardsasords]{lem:CardinalsBehaveRight}.
	\item إذا كان $\cardeq{\omega}{A}$، فإن $\card{A} = \omega$، بحسب
	\olref{omegaisacardinal}.
\end{enumerate}
إذن، في جميع الحالات، $\card{A} \leq \omega$.
\end{proof}
\noindent
للعدد $\omega$ بالفعل مكانة خاصة. فعلى الرغم من وجود أعداد ترتيبية كثيرة
قابلة للعد:

\begin{cor}
$\omega$ هو الكاردينال غير المتناهي الوحيد الذي هو !!{enumerable}.
\end{cor}

\begin{proof}
ليكن $\cardfont{a}$ كاردينالاً غير متناهٍ هو !!a{enumerable}. بما أن
$\cardfont{a}$ غير متناهٍ، فإن $\omega \leq \cardfont{a}$. وبما أن
$\cardfont{a}$ كاردينال هو !!a{enumerable}، فإن
$\cardfont{a} = \card{\cardfont{a}} \leq \omega$. ومن ثم
$\cardfont{a} = \omega$ بحسب قانون التثليث.
\end{proof}

وبالطبع توجد كاردينالات لا متناهية الكثرة. لذا قد نسأل:
\emph{كم كاردينالاً يوجد؟} تبين النتائج الآتية أننا قد نرغب في إعادة
النظر في هذا السؤال.

\begin{prop}\ollabel{unioncardinalscardinal}
إذا كان كل عنصر من $X$ كاردينالاً، فإن $\bigcup X$ كاردينال.
\end{prop}

\begin{proof}
يسهل التحقق من أن $\bigcup X$ عدد ترتيبي. وليكن
$\alpha \in \bigcup X$ عدداً ترتيبياً؛ عندئذ
$\alpha \in \cardfont{b} \in X$ لبعض الكاردينال $\cardfont{b}$. ولأن
$\cardfont{b}$ كاردينال، فإن $\cardless{\alpha}{\cardfont{b}}$. ولأن
$\cardfont{b} \subseteq \bigcup X$، فلدينا
$\cardle{\cardfont{b}}{\bigcup X}$، ومن ثم
$\cardneq{\alpha}{\bigcup X}$. وبتعميم ذلك يكون $\bigcup X$
كاردينالاً.
\end{proof}

\begin{thm}\ollabel{lem:NoLargestCardinal}
لا يوجد كاردينال أكبر.
\end{thm}

\begin{proof}
لكل كاردينال $\cardfont{a}$، تستلزم مبرهنة كانتور
(\olref[sfr][siz][car]{thm:cantor}) و
\olref[cardinals][cardsasords]{lem:CardinalsExist} أن
$\cardfont{a} < \card{\Pow{\cardfont{a}}}$.
\end{proof}

\begin{thm}
لا توجد مجموعة جميع الكاردينالات.
\end{thm}

\begin{proof}
للبرهان بالخلف، افترض أن
$C = \Setabs{\cardfont{a}}{\cardfont{a} \text{ كاردينال}}$. عندئذ
$\bigcup C$ كاردينال بحسب \olref{unioncardinalscardinal}، ولذلك يوجد،
بحسب \olref{lem:NoLargestCardinal}، كاردينال
$\cardfont{b} > \bigcup C$. وبحسب التعريف $\cardfont{b} \in C$، ومن
ثم $\cardfont{b} \subseteq \bigcup{C}$، وبالتالي
$\cardfont{b} \leq \bigcup C$، وهذا تناقض.
\end{proof}

ينبغي أن تقارن هذا بكل من مفارقة راسل ومفارقة بورالي--فورتي.

\end{document}
